\documentclass[11pt,a4paper]{article}

%----------------------------------------------------------------------------------------
%	PACKAGES
%----------------------------------------------------------------------------------------
\usepackage[left=0.8in,top=0.8in,right=0.8in,bottom=0.8in]{geometry} 
\usepackage[utf8]{inputenc}
\usepackage[T1]{fontenc}
\usepackage[default]{sourcesanspro} 
\usepackage{xcolor}      
\usepackage{titlesec}    
\usepackage{enumitem}    
\usepackage{parskip}     

%----------------------------------------------------------------------------------------
%	COLORS & STYLES
%----------------------------------------------------------------------------------------
\definecolor{titleblue}{HTML}{00199e} 
\definecolor{subtitleblue}{HTML}{2ec1e0} 
\definecolor{darktext}{HTML}{222222} 

\color{darktext}
\linespread{0.9} 
\setlength{\parskip}{0.1em} 

\titleformat{\section}{\Large\bfseries\color{titleblue}}{}{0em}{} 
\titlespacing*{\section}{0pt}{0.7em}{0.15em}

\newcommand{\blueitem}[1]{\textcolor{subtitleblue}{\textbf{#1}}}

\setlist[itemize]{label=\textbullet, leftmargin=*, noitemsep, topsep=0pt, parsep=0pt}

%----------------------------------------------------------------------------------------
%	DOCUMENT CONTENT
%----------------------------------------------------------------------------------------
\begin{document}

%--- HEADER ---
{\Huge \textbf{\textcolor{titleblue}{Lyara Nascimento}}} \vspace{0.2em}

Rio de Janeiro, Brasil \\
Celular: +55 21 971120503 \\
Email: lyaranascimento12@@gmail.com \\
LinkedIn: linkedin.com/in/lyaranascimento \\

\vspace{0.4em}

%--- PERSONAL PROFILE ---
\section*{Perfil Pessoal}
Sou cientista de dados e mestranda em Astronomia no Observatório do Valongo (UFRJ), com dupla graduação em Astronomia (UFRJ) e Tecnologia em Ciência de Dados (UNINTER). Tenho sólida experiência no desenvolvimento de pipelines computacionais, modelagem estatística e aplicação de algoritmos de Machine Learning para a análise de grandes volumes de dados astronômicos. Meu objetivo é continuar impulsionando a astrofísica computacional, unindo o rigor científico da pesquisa de sistemas planetários à inovação tecnológica da ciência de dados.

%--- EDUCATION ---
\section*{Educação}

\textbf{2028 -- Presente: Mestranda em Astronomia (Ênfase em Astrofísica Computacional e Ciência de Dados)}\\
Observatório do Valongo, Universidade Federal do Rio de Janeiro (UFRJ), Rio de Janeiro, RJ\\

\textbf{2023 -- 2028: Graduação em Astronomia}\\
Universidade Federal do Rio de Janeiro (UFRJ), Rio de Janeiro, RJ\\

\textbf{2025 -- 2028: Tecnologia em Ciência de Dados}\\
Centro Universitário Internacional (UNINTER), Rio de Janeiro, RJ\\

\textbf{2019 -- 2022: Ensino Médio}\\
Colégio Pedro II, Rio de Janeiro, RJ

%--- EXPERIENCE ---
\section*{Experiência Profissional}

\textbf{2028 -- Presente: Pesquisadora de Mestrado (Bolsista CNPq)}\\
Observatório do Valongo -- UFRJ
\begin{itemize}
    \item Desenvolvimento de projeto de pesquisa voltado à análise computacional de dados astronômicos.
    \item Tratamento e processamento de grandes volumes de dados de catálogos virtuais.
\end{itemize}

\textbf{2028 -- Presente: Cientista de Dados Júnior}\\
AstroData Analytics S.A.
\begin{itemize}
    \item Atuação na modelagem estatística e análise exploratória de grandes bancos de dados.
    \item Criação de dashboards de visualização e relatórios técnicos automatizados.
\end{itemize}

\textbf{2020 -- 2022: Estagiária de Iniciação Científica Jr. (Bolsista FAPERJ)}\\
Instituto de Pesquisas Jardim Botânico do Rio de Janeiro (IP/JBRJ)

\vspace{0.3em}


\end{itemize}

%--- PUBLICATIONS ---
\section*{Publicações Acadêmicas}

\textbf{Trabalho de Conclusão de Curso (TCC):}\\
NASCIMENTO, L. L. \textbf{Mineração de Dados no Catálogo TESS: Classificação Automatizada de Candidatos a Exoplanetas usando Machine Learning}. 2028. Graduação em Astronomia -- UFRJ.\\

\textbf{Iniciação Científica (Graduação):}\\
Pesquisa voltada à democratização do acesso à ciência, criando estratégias de divulgação astronômica baseadas em dados abertos para comunidades vulneráveis.

\blueitem{Contribuições científicas de mulheres botânicas na cidade do Rio de Janeiro desde o período colonial: um olhar sobre a história e a botânica.} \\
\textbf{Biologia e Meio ambiente: princípios, conceitos e práticas} Aya Editora, 2023. p. 42. 
\textit{\textbf{NASCIMENTO, L. L;} PANTOJA, S. C. S.}

\vspace{0.2em}


%--- PROJECTS ---

%--- SKILLS ---
\section*{Habilidades}
Python, Pacote Office \\
Inglês (B2), Portuguẽs (Fluente), Francês (A1)

%--- REFERENCES ---

\end{document}
